% !TEX program = xelatex
% !LW recipe = xelatex (2 passes)
% LuaLaTeX fails here: its font cache keeps the Titillium path relative to the slides folder.
\documentclass[11pt, a4paper]{article}

\usepackage[indonesian]{babel}
\usepackage{fontspec}
\setmainfont{TitilliumWeb}[
  Path = ../slides/theme/fonts/,
  Extension = .ttf,
  UprightFont = *-Regular,
  BoldFont = *-Bold,
  ItalicFont = *-Italic,
  BoldItalicFont = *-BoldItalic
]
\usepackage[margin=2.5cm]{geometry}
\usepackage{xcolor}
\usepackage{graphicx}
\usepackage{array}
\usepackage{booktabs}
\usepackage{tabularx}
\usepackage{enumitem}
\usepackage{titlesec}
\usepackage{tcolorbox}
\usepackage{caption}
% Keeps every table and figure inside its own section.
\usepackage[section]{placeins}
\usepackage{tikz}
\usetikzlibrary{arrows.meta, positioning, calc, decorations.pathreplacing}
\usepackage{hyperref}

\definecolor{praditagreen}{RGB}{37, 113, 71}

\hypersetup{
  colorlinks = true,
  linkcolor = praditagreen,
  citecolor = praditagreen,
  urlcolor = praditagreen,
  pdftitle = {Tugas Agentic Workflow},
  pdfauthor = {Alfa Yohannis}
}

\titleformat{\section}{\Large\bfseries\color{praditagreen}}{\thesection}{0.6em}{}
\titleformat{\subsection}{\large\bfseries\color{praditagreen}}{\thesubsection}{0.6em}{}

\captionsetup{font=small, labelfont=bf}
\setlist{itemsep=2pt, topsep=4pt}
\setlength{\parindent}{0pt}
\setlength{\parskip}{6pt}
\renewcommand{\arraystretch}{1.25}
\frenchspacing
\urlstyle{same}

% Let tables share a page with text instead of sitting alone on a float page.
\renewcommand{\topfraction}{0.9}
\renewcommand{\bottomfraction}{0.8}
\renewcommand{\textfraction}{0.07}
\renewcommand{\floatpagefraction}{0.85}

% Left-aligned columns, so narrow cells do not stretch the spaces between words.
\newcolumntype{Y}{>{\raggedright\arraybackslash}X}
\newcolumntype{L}[1]{>{\raggedright\arraybackslash}p{#1}}

% English terms that sound unnatural in Indonesian keep their English form, in italics.
\newcommand{\eng}[1]{\textit{#1}}

\begin{document}

\renewcommand{\refname}{Daftar Pustaka}

\begin{center}
  \includegraphics[height=1.8cm]{../slides/theme/classFiles/logo_universitas.png}\par
  \vspace{12pt}
  {\large\bfseries\color{praditagreen} Program Studi Magister Teknologi Informasi}\par
  \vspace{6pt}
  {\LARGE\bfseries Tugas: Penerapan \eng{Agentic Workflow} di Tempat Kerja}\par
  \vspace{6pt}
  Dosen: Alfa Yohannis
\end{center}

\begin{tcolorbox}[colback=praditagreen!5, colframe=praditagreen, boxrule=0.6pt, arc=2pt]
  \begin{tabularx}{\linewidth}{@{}L{3.6cm}Y@{}}
    \textbf{Bentuk tugas}    & Individu \\
    \textbf{Hasil akhir}     & Artikel ilmiah dan prototipe \eng{agentic workflow} \\
    \textbf{Panjang artikel} & Sekitar 6.000 sampai 7.500 kata \\
    \textbf{Metodologi}      & \eng{Design Science Research} (DSR) dalam bentuk studi kasus \\
    \textbf{Tenggat}         & Ditentukan dosen \\
  \end{tabularx}
\end{tcolorbox}

\section{Gambaran Tugas}

Tugas berupa penelitian tentang penerapan \eng{agentic workflow} pada satu proses kerja nyata.
Objek penelitian adalah proses kerja di tempat mahasiswa bekerja.
Tempat kerja boleh berasal dari semua bidang, misalnya sekolah, bank, perusahaan perangkat lunak, atau kantor akuntan.

Tiga pertanyaan penelitian (\eng{research question}, RQ) berlaku sama untuk semua bidang.
Perbedaan hanya terletak pada proses kerja yang dipilih.
Hasil akhir berupa satu artikel ilmiah dan satu prototipe sederhana.

Setelah menyelesaikan tugas, mahasiswa mampu:
\begin{enumerate}
  \item menilai proses kerja yang cocok untuk \eng{agentic workflow};
  \item merancang \eng{agentic workflow} yang tetap diawasi manusia;
  \item mengukur manfaat dan risiko \eng{agentic workflow} berdasarkan data.
\end{enumerate}

\section{Pengertian \eng{Agentic Workflow}}

\eng{Agentic workflow} adalah alur kerja yang melibatkan agen kecerdasan buatan (\eng{artificial intelligence}, AI) untuk mengerjakan sebagian langkah kerja.
Agen AI adalah program berbasis \eng{large language model} (LLM) yang mampu merencanakan langkah, memakai alat, dan memeriksa hasil kerja.
Alat dapat berupa basis data, \eng{email}, \eng{spreadsheet}, atau aplikasi internal.
Manusia tetap mengambil keputusan penting dan memeriksa hasil pada titik tertentu.

Tabel~\ref{tab:beda} membandingkan \eng{agentic workflow} dengan otomatisasi biasa, misalnya skrip dan \eng{robotic process automation} (RPA).
Tabel~\ref{tab:contoh} memberi contoh penerapan di empat bidang.
Contoh hanya berfungsi sebagai gambaran.
Proses yang diteliti tetap berasal dari tempat mahasiswa bekerja.

\begin{table}[htbp]
  \centering
  \caption{Perbedaan otomatisasi biasa dan \eng{agentic workflow}}
  \label{tab:beda}
  \begin{tabularx}{\linewidth}{@{}L{3.4cm}YY@{}}
    \toprule
    \textbf{Aspek} & \textbf{Otomatisasi biasa (skrip, RPA)} & \textbf{\eng{Agentic workflow}} \\
    \midrule
    Cara kerja             & Mengikuti aturan tetap dari programmer & Menyusun langkah secara mandiri berdasarkan tujuan \\
    Jenis masukan          & Data terstruktur dengan format tetap   & Data terstruktur, teks bebas, dokumen, dan \eng{email} \\
    Saat muncul kasus baru & Berhenti atau gagal                    & Mencoba cara lain atau bertanya kepada manusia \\
    Risiko utama           & Aturan tidak lengkap                   & Hasil salah yang tampak meyakinkan \\
    \bottomrule
  \end{tabularx}
\end{table}

\begin{table}[htbp]
  \centering
  \caption{Contoh \eng{agentic workflow} di empat bidang}
  \label{tab:contoh}
  \begin{tabularx}{\linewidth}{@{}L{2.4cm}L{3cm}YY@{}}
    \toprule
    \textbf{Bidang} & \textbf{Contoh proses} & \textbf{Peran agen AI} & \textbf{Peran manusia} \\
    \midrule
    Sekolah & Penyusunan soal ujian
      & Menyusun draf soal dan kunci jawaban dari silabus
      & Guru memeriksa dan menyetujui soal \\
    Bank & Pemeriksaan dokumen pengajuan kredit
      & Memeriksa kelengkapan dokumen dan menyiapkan \eng{email} untuk nasabah
      & Analis kredit memutuskan kelanjutan pengajuan \\
    Perusahaan perangkat lunak & Penanganan laporan \eng{bug}
      & Mereproduksi \eng{bug}, mengusulkan perbaikan, dan menjalankan tes
      & Pengembang meninjau dan menyetujui perbaikan \\
    Kantor akuntan & Rekonsiliasi bank
      & Mencocokkan transaksi dan menandai selisih
      & Akuntan memeriksa selisih dan menyetujui jurnal \\
    \bottomrule
  \end{tabularx}
\end{table}

\section{Pertanyaan Penelitian}

Ketiga RQ mengikuti urutan kerja: memilih proses, merancang solusi, lalu mengukur hasil.
Kata ``organisasi'' boleh diganti dengan jenis tempat kerja, misalnya ``sekolah'' atau ``bank''.
Inti pertanyaan tidak boleh berubah.

\begin{table}[htbp]
  \centering
  \caption{Tiga pertanyaan penelitian}
  \label{tab:rq}
  \begin{tabularx}{\linewidth}{@{}L{1.1cm}YL{4.6cm}@{}}
    \toprule
    \textbf{RQ} & \textbf{Pertanyaan} & \textbf{Bentuk jawaban} \\
    \midrule
    RQ1 & Proses kerja mana di organisasi yang cocok untuk \eng{agentic workflow}, dan ciri apa yang menentukan kecocokan proses?
        & Daftar proses kandidat dengan skor kelayakan, serta satu proses terpilih \\
    RQ2 & Bagaimana rancangan \eng{agentic workflow} untuk proses terpilih, meliputi peran agen, alat, batas wewenang agen, dan titik pemeriksaan oleh manusia?
        & Diagram alur kerja baru dan prototipe yang dapat dijalankan \\
    RQ3 & Seberapa besar pengaruh \eng{agentic workflow} terhadap waktu proses, kualitas hasil, dan beban kerja manusia, serta risiko apa yang muncul?
        & Angka perbandingan cara lama dan cara baru, pendapat pengguna, dan daftar risiko \\
    \bottomrule
  \end{tabularx}
\end{table}

\section{Metode untuk Setiap RQ}

Setiap RQ dijawab dengan tiga metode.
Ketiga metode saling melengkapi agar jawaban RQ tidak bergantung pada satu sumber data.
Setiap metode diberi kode, misalnya M1.2 untuk metode kedua pada RQ1.
Kode yang sama dipakai lagi pada Gambar~\ref{fig:dsr}.

\subsection{Metode untuk RQ1: Memilih Proses}

Mahasiswa menilai minimal tiga proses kandidat, lalu memilih satu proses untuk tahap berikutnya.

\begin{table}[htbp]
  \centering
  \caption{Metode untuk RQ1}
  \label{tab:metode-rq1}
  \begin{tabularx}{\linewidth}{@{}L{1cm}L{2.6cm}YL{3.2cm}@{}}
    \toprule
    \textbf{Kode} & \textbf{Metode} & \textbf{Cara pelaksanaan} & \textbf{Hasil} \\
    \midrule
    M1.1 & Wawancara
         & Wawancarai 3 sampai 5 orang yang menjalankan atau memimpin proses. Tanyakan langkah kerja, masalah utama, dan waktu yang dihabiskan.
         & Catatan wawancara dan daftar masalah \\
    M1.2 & Pemetaan proses saat ini
         & Gambarkan alur kerja saat ini dalam diagram alir atau \eng{Business Process Model and Notation} (BPMN). Catat pelaku dan perkiraan waktu setiap langkah. Sumber informasi berupa pengamatan langsung, standar operasional prosedur (SOP), atau \eng{log} sistem.
         & Diagram proses saat ini \\
    M1.3 & Penilaian kelayakan
         & Beri skor 1 sampai 5 untuk setiap proses kandidat pada lima kriteria di Tabel~\ref{tab:kriteria}. Jumlahkan skor, lalu pilih proses dengan skor tertinggi.
         & Tabel skor dan satu proses terpilih \\
    \bottomrule
  \end{tabularx}
\end{table}

\begin{table}[htbp]
  \centering
  \caption{Kriteria kelayakan proses (skor 5 berarti paling cocok)}
  \label{tab:kriteria}
  \begin{tabularx}{\linewidth}{@{}L{4.4cm}Y@{}}
    \toprule
    \textbf{Kriteria} & \textbf{Pertanyaan panduan} \\
    \midrule
    Frekuensi                 & Seberapa sering proses terjadi? \\
    Keteraturan langkah       & Apakah langkah kerja cukup jelas dan berulang? \\
    Ketersediaan data digital & Apakah data dan dokumen sudah tersedia dalam bentuk digital? \\
    Kemudahan koreksi         & Apakah kesalahan agen mudah dikenali dan diperbaiki? \\
    Beban kerja manusia       & Seberapa banyak waktu staf yang tersita oleh proses? \\
    \bottomrule
  \end{tabularx}
\end{table}

Kriteria boleh ditambah sesuai kebutuhan bidang.
Arah skor setiap kriteria harus sama, yaitu skor tinggi berarti makin cocok.

\subsection{Metode untuk RQ2: Merancang Solusi}

Setiap rancangan wajib memiliki minimal satu titik persetujuan manusia.
Makin tinggi risiko sebuah langkah, makin rendah tingkat wewenang agen pada langkah tersebut.

\begin{table}[htbp]
  \centering
  \caption{Metode untuk RQ2}
  \label{tab:metode-rq2}
  \begin{tabularx}{\linewidth}{@{}L{1cm}L{2.6cm}YL{3.2cm}@{}}
    \toprule
    \textbf{Kode} & \textbf{Metode} & \textbf{Cara pelaksanaan} & \textbf{Hasil} \\
    \midrule
    M2.1 & Perancangan alur kerja baru
         & Gambarkan alur kerja baru dengan notasi yang sama seperti M1.2. Tandai langkah oleh agen, langkah oleh manusia, dan titik persetujuan manusia. Tentukan tingkat wewenang agen untuk setiap langkah berdasarkan Tabel~\ref{tab:wewenang}.
         & Diagram alur kerja baru dan tabel wewenang agen \\
    M2.2 & Pembuatan prototipe
         & Bangun prototipe untuk proses inti dengan data contoh. Alat bebas dipilih, misalnya n8n, LangGraph, CrewAI, atau Claude Agent SDK (\eng{software development kit}). Catat model AI dan \eng{prompt} yang dipakai.
         & Prototipe yang dapat didemokan, beserta kode atau tangkapan layar \\
    M2.3 & Tinjauan pakar
         & Minta 2 sampai 3 orang berpengalaman di tempat kerja menilai rancangan dan demo prototipe. Pakai daftar pertanyaan yang sama untuk setiap pakar. Perbaiki rancangan berdasarkan masukan pakar.
         & Ringkasan masukan pakar dan daftar perbaikan \\
    \bottomrule
  \end{tabularx}
\end{table}

\begin{table}[htbp]
  \centering
  \caption{Tiga tingkat wewenang agen}
  \label{tab:wewenang}
  \begin{tabularx}{\linewidth}{@{}L{2.6cm}YYY@{}}
    \toprule
    \textbf{Tingkat} & \textbf{Peran agen} & \textbf{Peran manusia} & \textbf{Cocok untuk} \\
    \midrule
    1. Usulan      & Memberi usulan atau draf                      & Memutuskan dan mengerjakan       & Langkah berisiko tinggi \\
    2. Persetujuan & Mengerjakan langkah                           & Menyetujui hasil sebelum dipakai & Langkah berisiko sedang \\
    3. Mandiri     & Mengerjakan langkah dan langsung memakai hasil & Memeriksa hasil secara berkala  & Langkah berisiko rendah dan mudah dikoreksi \\
    \bottomrule
  \end{tabularx}
\end{table}

\subsection{Metode untuk RQ3: Mengukur Hasil}

Semua hasil dilaporkan apa adanya, termasuk hasil yang tidak sesuai harapan.
Ukuran tambahan boleh dipakai sesuai kebutuhan bidang.

\begin{table}[htbp]
  \centering
  \caption{Metode untuk RQ3}
  \label{tab:metode-rq3}
  \begin{tabularx}{\linewidth}{@{}L{1cm}L{2.6cm}YL{3.2cm}@{}}
    \toprule
    \textbf{Kode} & \textbf{Metode} & \textbf{Cara pelaksanaan} & \textbf{Hasil} \\
    \midrule
    M3.1 & Uji perbandingan
         & Siapkan minimal 20 kasus nyata yang sudah dianonimkan. Kerjakan setiap kasus dengan cara lama dan dengan \eng{agentic workflow}. Bandingkan hasil dengan ukuran pada Tabel~\ref{tab:ukuran}.
         & Tabel perbandingan cara lama dan cara baru \\
    M3.2 & Kuesioner dan wawancara pengguna
         & Minta minimal 5 pengguna mencoba prototipe. Pengguna mengisi kuesioner \eng{System Usability Scale} (SUS)~\cite{brooke1996} dan beberapa pertanyaan tentang kepercayaan pada hasil agen. SUS berisi 10 pernyataan dengan skala 1 sampai 5 dan menghasilkan skor 0 sampai 100. Lanjutkan dengan wawancara singkat.
         & Skor SUS, skor kepercayaan, dan kutipan pendapat pengguna \\
    M3.3 & Analisis kesalahan dan risiko
         & Catat setiap kesalahan agen selama uji, lalu kelompokkan menurut jenis kesalahan. Uji juga kasus sulit, misalnya data tidak lengkap atau perintah yang menyesatkan. Gunakan \eng{OWASP Top 10 for LLM Applications}~\cite{owasp2025} sebagai daftar periksa. OWASP adalah \eng{Open Worldwide Application Security Project}.
         & Daftar risiko dan cara pengendalian \\
    \bottomrule
  \end{tabularx}
\end{table}

\begin{table}[htbp]
  \centering
  \caption{Ukuran utama untuk uji perbandingan}
  \label{tab:ukuran}
  \begin{tabularx}{\linewidth}{@{}L{4.4cm}Y@{}}
    \toprule
    \textbf{Ukuran} & \textbf{Cara menghitung} \\
    \midrule
    Waktu proses        & Menit sejak kasus masuk sampai selesai \\
    Waktu kerja manusia & Menit kerja staf untuk setiap kasus \\
    Tingkat kesalahan   & Persentase kasus dengan hasil salah \\
    Biaya per kasus     & Biaya layanan AI ditambah biaya waktu staf \\
    \bottomrule
  \end{tabularx}
\end{table}

\section{Metodologi yang Direkomendasikan}

Metodologi yang direkomendasikan adalah \eng{Design Science Research} (DSR) dalam bentuk studi kasus di satu tempat kerja.
DSR adalah metode penelitian yang menghasilkan solusi bagi masalah nyata, lalu menguji solusi secara terukur.
Tahapan DSR mengacu pada Peffers, Tuunanen, Rothenberger, dan Chatterjee~\cite{peffers2007}.

Alasan pemilihan DSR:
\begin{enumerate}
  \item Tugas menuntut pembuatan solusi, bukan hanya pengamatan.
  \item Tahapan DSR jelas dan sudah lazim dalam penelitian sistem informasi.
  \item Tahapan yang sama berlaku untuk semua bidang tempat kerja.
\end{enumerate}

Gambar~\ref{fig:dsr} menunjukkan enam tahap DSR dan hubungan setiap tahap dengan RQ.
Bila hasil evaluasi belum memenuhi target, penelitian kembali ke tahap desain.
Setiap putaran perbaikan dicatat dalam artikel.

Data penelitian terdiri atas dua jenis.
Data kualitatif berasal dari wawancara dan tinjauan pakar.
Data kuantitatif berasal dari skor kelayakan, uji perbandingan, dan kuesioner.

\begin{figure}[htbp]
  \centering
  \begin{tikzpicture}[
      tahap/.style={draw=praditagreen, thick, rounded corners=4pt, text width=8.6cm,
        minimum height=1.25cm, inner sep=6pt, align=left, font=\small},
      panah/.style={-{Stealth[length=2.4mm]}, thick, draw=black!70},
      kurung/.style={decorate, decoration={brace, amplitude=6pt}, thick, draw=black!70},
      label/.style={right=10pt, align=left, font=\small},
      node distance=5mm
    ]
    \node[tahap] (t1) {\textbf{1. Identifikasi masalah} (M1.1, M1.2)\\{\footnotesize Wawancara dan pemetaan proses saat ini}};
    \node[tahap, below=of t1] (t2) {\textbf{2. Penentuan tujuan solusi} (M1.3)\\{\footnotesize Penilaian kelayakan dan target perbaikan}};
    \node[tahap, below=of t2] (t3) {\textbf{3. Desain dan pengembangan} (M2.1, M2.2)\\{\footnotesize Alur kerja baru dan prototipe}};
    \node[tahap, below=of t3] (t4) {\textbf{4. Demonstrasi} (M2.3)\\{\footnotesize Demo prototipe kepada pakar}};
    \node[tahap, below=of t4] (t5) {\textbf{5. Evaluasi} (M3.1, M3.2, M3.3)\\{\footnotesize Uji perbandingan, kuesioner, dan analisis risiko}};
    \node[tahap, below=of t5] (t6) {\textbf{6. Komunikasi}\\{\footnotesize Penulisan artikel ilmiah}};

    \foreach \a/\b in {t1/t2, t2/t3, t3/t4, t4/t5, t5/t6}
      \draw[panah] (\a) -- (\b);

    % Loop back from evaluation to design.
    \draw[panah] (t5.west) -- ++(-0.7, 0) |- (t3.west);
    \node[font=\footnotesize, align=right, anchor=east]
      at ($(t4.west) + (-0.85, 0)$) {ulangi bila\\hasil belum\\memenuhi\\target};

    % Which research question each group of stages answers.
    \draw[kurung] ($(t1.north east) + (0.25, 0)$) -- ($(t2.south east) + (0.25, 0)$)
      node[midway, label] {\textbf{RQ1}\\Proses mana\\yang cocok?};
    \draw[kurung] ($(t3.north east) + (0.25, 0)$) -- ($(t4.south east) + (0.25, 0)$)
      node[midway, label] {\textbf{RQ2}\\Rancangan\\seperti apa?};
    \draw[kurung] ($(t5.north east) + (0.25, 0)$) -- ($(t5.south east) + (0.25, 0)$)
      node[midway, label] {\textbf{RQ3}\\Apa manfaat\\dan risiko?};
    \draw[kurung] ($(t6.north east) + (0.25, 0)$) -- ($(t6.south east) + (0.25, 0)$)
      node[midway, label] {\textbf{Artikel}};
  \end{tikzpicture}
  \caption{Enam tahap DSR~\cite{peffers2007} dan hubungan dengan RQ}
  \label{fig:dsr}
\end{figure}

\section{Struktur Artikel}

Struktur artikel mengikuti pola artikel DSR dari Gregor dan Hevner~\cite{gregor2013}.
Penyesuaian dilakukan agar setiap RQ memiliki satu bagian hasil.
Panjang artikel sekitar 6.000 sampai 7.500 kata, tidak termasuk daftar pustaka dan lampiran.

\begin{table}[htbp]
  \centering
  \caption{Struktur artikel yang direkomendasikan}
  \label{tab:struktur}
  \small
  \begin{tabularx}{\linewidth}{@{}L{1.3cm}L{3.6cm}YL{3.1cm}@{}}
    \toprule
    \textbf{Nomor} & \textbf{Bagian} & \textbf{Isi utama} & \textbf{Panjang (kata)} \\
    \midrule
       & Judul, abstrak, kata kunci & Masalah, metode, dan hasil utama. Tiga sampai lima kata kunci. & 150 sampai 250 \\
    1  & Pendahuluan & Latar belakang, masalah di tempat kerja, tiga RQ, dan kontribusi penelitian & 600 sampai 800 \\
    2  & Tinjauan pustaka & Konsep \eng{agentic workflow}, penelitian terdahulu, dan perbedaan dengan otomatisasi biasa & 800 sampai 1.000 \\
    3  & Metodologi & Tahapan DSR, profil tempat kerja tanpa nama asli, cara pengumpulan dan analisis data, serta etika & 600 sampai 800 \\
    4  & Analisis proses kerja (RQ1) & Hasil wawancara, diagram proses saat ini, tabel skor, dan proses terpilih & 800 sampai 1.000 \\
    5  & Rancangan \eng{agentic workflow} (RQ2) & Diagram alur kerja baru, peran agen, alat, tabel wewenang, prototipe, dan masukan pakar & 1.000 sampai 1.200 \\
    6  & Evaluasi (RQ3) & Cara uji, tabel perbandingan, hasil kuesioner, dan daftar risiko & 1.000 sampai 1.200 \\
    7  & Pembahasan & Jawaban singkat setiap RQ, manfaat bagi tempat kerja, dan keterbatasan penelitian & 600 sampai 800 \\
    8  & Kesimpulan dan saran & Temuan utama dan saran penelitian lanjutan & 300 sampai 400 \\
       & Daftar pustaka & Minimal 15 sumber dengan gaya \eng{American Psychological Association} (APA) edisi ke-7 & \\
       & Lampiran & Pedoman wawancara, kuesioner, contoh \eng{prompt}, dan contoh data uji yang sudah dianonimkan & \\
    \bottomrule
  \end{tabularx}
\end{table}

Panduan penulisan:
\begin{itemize}
  \item Setiap bagian hasil dibuka dengan jawaban singkat atas RQ, lalu diikuti bukti.
  \item Setiap gambar dan tabel diberi nomor dan dirujuk dalam teks.
  \item Profil tempat kerja ditulis tanpa nama asli, misalnya ``sebuah bank swasta di Jakarta''.
\end{itemize}

\section{Etika dan Kerahasiaan Data}

Penelitian memakai data dari tempat kerja nyata, sehingga aturan berikut wajib dipatuhi:
\begin{itemize}
  \item Izin tertulis dari atasan atau pihak berwenang diperoleh sebelum pengambilan data.
  \item Setiap responden memberi persetujuan sebelum wawancara dan pengisian kuesioner.
  \item Nama organisasi, nama orang, dan data pribadi wajib dianonimkan, misalnya nomor rekening, nilai siswa, atau data gaji.
  \item Data rahasia tidak boleh dimasukkan ke layanan AI publik. Gunakan data contoh atau data yang sudah disamarkan.
  \item Prototipe tidak boleh terhubung ke sistem produksi tanpa izin.
  \item Pengolahan data pribadi mengikuti Undang-Undang Nomor 27 Tahun 2022 tentang Pelindungan Data Pribadi~\cite{uupdp2022}.
  \item Penggunaan alat AI dalam penulisan artikel wajib dinyatakan pada bagian metodologi.
\end{itemize}

\section{Yang Dikumpulkan dan Penilaian}

Mahasiswa mengumpulkan tiga hal:
\begin{enumerate}
  \item artikel dalam bentuk berkas \eng{Portable Document Format} (PDF);
  \item bukti prototipe, berupa tautan repositori kode atau video demo berdurasi 3 sampai 5 menit;
  \item lampiran instrumen penelitian dan data uji yang sudah dianonimkan.
\end{enumerate}

Tabel~\ref{tab:nilai} memuat kriteria penilaian.

\begin{table}[htbp]
  \centering
  \caption{Kriteria penilaian}
  \label{tab:nilai}
  \begin{tabularx}{\linewidth}{@{}L{4.4cm}YL{1.4cm}@{}}
    \toprule
    \textbf{Kriteria} & \textbf{Yang dinilai} & \textbf{Bobot} \\
    \midrule
    Analisis proses kerja (RQ1)   & Wawancara memadai, diagram proses jelas, dan skor kelayakan masuk akal & 20\% \\
    Rancangan dan prototipe (RQ2) & Peran agen jelas, titik pemeriksaan manusia tepat, dan prototipe berjalan & 25\% \\
    Evaluasi (RQ3)                & Data uji cukup, hasil dilaporkan apa adanya, dan risiko dibahas & 25\% \\
    Pembahasan dan kesimpulan     & Jawaban RQ jelas dan keterbatasan penelitian diakui & 15\% \\
    Kualitas penulisan            & Struktur rapi, bahasa jelas, dan rujukan benar & 15\% \\
    \midrule
    \textbf{Total}                & & \textbf{100\%} \\
    \bottomrule
  \end{tabularx}
\end{table}

\begin{thebibliography}{9}

\bibitem{brooke1996}
Brooke, J. (1996).
SUS: A quick and dirty usability scale.
Dalam P.~W. Jordan, B. Thomas, I.~L. McClelland, \& B. Weerdmeester (Editor),
\textit{Usability evaluation in industry} (halaman 189--194). Taylor \& Francis.
\url{https://www.routledge.com/products/9780429157011}

\bibitem{gregor2013}
Gregor, S., \& Hevner, A.~R. (2013).
Positioning and presenting design science research for maximum impact.
\textit{MIS Quarterly}, \textit{37}(2), 337--355.
\url{https://doi.org/10.25300/MISQ/2013/37.2.01}

\bibitem{owasp2025}
OWASP GenAI Security Project. (2024).
\textit{OWASP Top 10 for LLM applications 2025}. OWASP Foundation.
\url{https://genai.owasp.org/llm-top-10/}

\bibitem{peffers2007}
Peffers, K., Tuunanen, T., Rothenberger, M.~A., \& Chatterjee, S. (2007).
A design science research methodology for information systems research.
\textit{Journal of Management Information Systems}, \textit{24}(3), 45--77.
\url{https://doi.org/10.2753/MIS0742-1222240302}

\bibitem{uupdp2022}
Republik Indonesia. (2022).
\textit{Undang-Undang Nomor 27 Tahun 2022 tentang Pelindungan Data Pribadi}.
\url{https://peraturan.bpk.go.id/Details/229798/uu-no-27-tahun-2022}

\end{thebibliography}

\end{document}
